



%\begin{table}[t]
%\centering
%\scalebox{0.7}{
%\begin{tabular}{lccccc|ccccc}
%\cline{7-11}
%\toprule[1.2pt]
% & \multicolumn{5}{c|}{Energy MAE (eV) $\downarrow$} & \multicolumn{5}{c}{EwT (\%) $\uparrow$} \\ 
%\cmidrule[0.6pt]{2-11}
%Methods & ID & OOD Ads & OOD Cat & OOD Both & Average & ID & OOD Ads & OOD Cat & OOD Both & Average \\
%\midrule[1.2pt]
%GNS~\cite{noisy_nodes} & 0.54 & 0.65 & 0.55 & 0.59 & 0.5825 & - & - & - & - & -\\
%GNS + Noisy Nodes~\cite{noisy_nodes} & 0.48 & 0.53 & 0.49 & 0.48 & 0.4950 & - & - & - & - & - \\

%Noisy Nodes~\cite{noisy_nodes} & 0.47 & 0.51 & 0.48 & 0.46 & 0.4800 & - & - & - & - & - \\
%Graphormer~\cite{graphormer_3d} & 0.4329 & 0.5850 & 0.4441 & 0.5299 & 0.4980 & - & - & - & - & -\\
%\midrule[0.6pt]
%Equiformer& 0.5093& 0.6285 & 0.5053 & 0.5556 & 0.5497 & 5.04 & 2.85 & 4.90 & 3.04 & \\
%Equiformer & 0.4222 & 0.5420 & 0.4231 & 0.4754 & 0.4657 & 7.23 & 3.77 & 7.13 & 4.10 & 5.56 \\
%\midrule[0.6pt]
%Equiformer + NAT & 0.4685 & 0.6071 & 0.4746 & 0.5345 & 0.5212 & 5.83 & 3.18 & 5.88 & 3.43 & \\
%\bottomrule[1.2pt]
%\end{tabular}
%}
%\vspace{2mm}
%\caption{\textbf{Results on OC20 IS2RE %\underline{validation} set when IS2RS node-level auxiliary task is adopted during training.} 
%``GNS'' denotes the 50-layer GNS trained without Noisy Nodes data augmentation, and ``Noisy Nodes'' denotes the 100-layer GNS trained with Noisy Nodes.}
%\vspace{-7mm}
%\label{tab:is2re_is2rs_validation}
%\end{table}



\begin{table}[t!]
\begin{adjustwidth}{0cm}{0cm}
\centering
\resizebox{\ifdim\width>\textwidth\textwidth\else\width\fi}{!}{
\begin{tabular}{lccccc|ccccc}
%\cline{7-11}
\toprule[1.2pt]
 & \multicolumn{5}{c|}{Energy MAE (eV) $\downarrow$} & \multicolumn{5}{c}{EwT (\%) $\uparrow$} \\ 
\cmidrule[0.6pt]{2-11}
Methods & ID & OOD Ads & OOD Cat & OOD Both & Average & ID & OOD Ads & OOD Cat & OOD Both & Average \\
\midrule[1.2pt]
CGCNN~\citep{cgcnn} & 0.6149 & 0.9155 & 0.6219 & 0.8511 & 0.7509 & 3.40 & 1.93 & 3.10 & 2.00 & 2.61 \\
SchNet~\citep{schnet} & 0.6387 & 0.7342 & 0.6616 & 0.7037 & 0.6846 & 2.96 & 2.33 & 2.94 & 2.21 & 2.61 \\
DimeNet++~\citep{dimenet_pp} & 0.5621 & 0.7252 & 0.5756 & 0.6613 & 0.6311 & 4.25 & 2.07 & 4.10 & 2.41 & 3.21 \\
PaiNN~\citep{painn} &  0.575 & 0.783 & 0.604 & 0.743 & 0.6763 & 3.46 & 1.97 & 3.46 & 2.28 & 2.79 \\
SpinConv~\citep{spinconv} & 0.5583 & 0.7230 & 0.5687 & 0.6738 & 0.6310 & 4.08 & 2.26 & 3.82 & 2.33 & 3.12 \\
SphereNet~\citep{spherenet} & 0.5625 & 0.7033 & 0.5708 & 0.6378 & 0.6186 & 4.47 & 2.29 & 4.09 & 2.41 & 3.32 \\

%EGNN& 0.5497& 0.6851& 0.5519& 0.6102& 0.5992& 4.99& 2.50& 4.71& 2.88& \\
SEGNN~\citep{segnn} & 0.5327 & 0.6921 & 0.5369 & 0.6790 & 0.6101 & \textbf{5.37} & \textbf{2.46} & \textbf{4.91} & 2.63 & \textbf{3.84} \\
\midrule[0.6pt]
%Equiformer& 0.5093& 0.6285 & 0.5053 & 0.5556 & 0.5497 & 5.04 & 2.85 & 4.90 & 3.04 & \\
Equiformer & \textbf{0.5037} & \textbf{0.6881} & \textbf{0.5213} & \textbf{0.6301} & \textbf{0.5858} & 5.14 & 2.41 & 4.67 & \textbf{2.69} & 3.73 \\
%\midrule[0.6pt]
%Equiformer + NAT & 0.4737 & 0.7245 & 0.4862 & 0.6493 & 0.5834 & 6.01 & 2.45 & 5.41 & 2.71 & \\
\bottomrule[1.2pt]

\end{tabular}
}
\vspace{-3mm}
\end{adjustwidth}
\caption{\textbf{Results on OC20 IS2RE \underline{testing} set.}}
\vspace{-3mm}
\label{tab:is2re_test}
\end{table}


\begin{table}[t!]
\begin{adjustwidth}{0cm}{0cm}
\centering
\resizebox{\ifdim\width>\textwidth\textwidth\else\width\fi}{!}{
\begin{tabular}{lccccc|ccccc}
%\cline{7-11}
\toprule[1.2pt]
 & \multicolumn{5}{c|}{Energy MAE (eV) $\downarrow$} & \multicolumn{5}{c}{EwT (\%) $\uparrow$} \\ 
\cmidrule[0.6pt]{2-11}
Methods & ID & OOD Ads & OOD Cat & OOD Both & Average & ID & OOD Ads & OOD Cat & OOD Both & Average \\
\midrule[1.2pt]
GNS~\citep{noisy_nodes} & 0.54 & 0.65 & 0.55 & 0.59 & 0.5825 & - & - & - & - & -\\
%GNS + Noisy Nodes~\cite{noisy_nodes} & 0.48 & 0.53 & 0.49 & 0.48 & 0.4950 & - & - & - & - & - \\

GNS + Noisy Nodes~\citep{noisy_nodes} & 0.47 & 0.51 & 0.48 & 0.46 & 0.4800 & - & - & - & - & - \\
Graphormer~\citep{graphormer_3d} & 0.4329 & 0.5850 & 0.4441 & 0.5299 & 0.4980 & - & - & - & - & -\\
\midrule[0.6pt]
%Equiformer& 0.5093& 0.6285 & 0.5053 & 0.5556 & 0.5497 & 5.04 & 2.85 & 4.90 & 3.04 & \\
Equiformer & 0.4222 & 0.5420 & 0.4231 & 0.4754 & 0.4657 & 7.23 & 3.77 & 7.13 & 4.10 & 5.56 \\
Equiformer + Noisy Nodes & \textbf{0.4156} &  \textbf{0.4976} & \textbf{0.4165} & \textbf{0.4344} & \textbf{0.4410} & \textbf{7.47} & \textbf{4.64} & \textbf{7.19} & \textbf{4.84} & \textbf{6.04} \\
%\midrule[0.6pt]
%Equiformer + NAT & 0.4685 & 0.6071 & 0.4746 & 0.5345 & 0.5212 & 5.83 & 3.18 & 5.88 & 3.43 & \\
\bottomrule[1.2pt]
\end{tabular}
}
\end{adjustwidth}
\vspace{-3mm}
%%%\caption{\textbf{Results on OC20 IS2RE \underline{validation} set when IS2RS node-level auxiliary task is adopted during training.} 
\caption{\textbf{Results on OC20 IS2RE \underline{validation} set when IS2RS is adopted during training.} 
%%%``GNS'' denotes the 50-layer GNS trained without Noisy Nodes data augmentation, and ``GNS + Noisy Nodes'' denotes the 100-layer GNS trained with Noisy Nodes.
%Compared to the main text, we add the result of ``Equiformer + Noisy Nodes'', which use data augmentation of interpolating between initial structure and relaxed struture and adding Gaussian noise as described by Noisy Nodes~\cite{noisy_nodes}.
%%%``Equiformer + Noisy Nodes'' uses data augmentation of interpolating between initial structure and relaxed structure and adding Gaussian noise as described by Noisy Nodes~\citep{noisy_nodes}.
}
\vspace{-3mm}
\label{tab:is2re_is2rs_validation}
\end{table}


\begin{table}[t!]
\begin{adjustwidth}{0cm}{0cm}
\centering
\resizebox{\ifdim\width>\textwidth\textwidth\else\width\fi}{!}{
\begin{tabular}{lccccc|cccccb{0.15cm}b{0.15cm}}
%\cline{7-11}
\toprule[1.2pt]
 & \multicolumn{5}{c|}{Energy MAE (eV) $\downarrow$} & \multicolumn{5}{c}{EwT (\%) $\uparrow$} & \multicolumn{2}{c}{Training time} \\ 
\cmidrule[0.6pt]{2-11}
Methods & ID & OOD Ads & OOD Cat & OOD Both & Average & ID & OOD Ads & OOD Cat & OOD Both & Average & \multicolumn{2}{c}{(GPU-days)} \\
\midrule[1.2pt]
GNS + Noisy Nodes~\citep{noisy_nodes} & 0.4219 & 0.5678 & 0.4366 & \textbf{0.4651} & 0.4728 & \textbf{9.12} & \textbf{4.25} & 8.01 & \textbf{4.64} & \textbf{6.5} & \multicolumn{1}{b{0.17cm}}{56} & \multicolumn{1}{l}{(TPU)} \\
%GNS + Noisy Nodes~\cite{noisy_nodes} & 0.48 & 0.53 & 0.49 & 0.48 & 0.4950 & - & - & - & - & - \\
%
%Noisy Nodes~\cite{noisy_nodes} & 0.47 & 0.51 & 0.48 & 0.46 & 0.4800 & - & - & - & - & - \\
Graphormer~\citep{graphormer_3d}$^\dagger$ & \textbf{0.3976} & 0.5719 & \textbf{0.4166} & 0.5029 & 0.4722 & 8.97 & 3.45 & \textbf{8.18} & 3.79 & 6.1 & \multicolumn{1}{b{0.17cm}}{372} & \multicolumn{1}{l}{(A100)} \\ %372 (A100)\\
\midrule[0.6pt]
%Equiformer& 0.5093& 0.6285 & 0.5053 & 0.5556 & 0.5497 & 5.04 & 2.85 & 4.90 & 3.04 & \\
Equiformer + Noisy Nodes & 0.4171 & \textbf{0.5479} & 0.4248 & 0.4741 & \textbf{0.4660} & 7.71& 3.70 & 7.15 & 4.07 & 5.66 & \multicolumn{1}{b{0.17cm}}{\textbf{24}} & \multicolumn{1}{l}{(A6000)} \\ %\multicolumn{1}{r}{24   (A6000)} \\
%+ Noisy Nodes & 0.4156 &  0.4976 & 0.4165 & 0.4344 & 0.4410 & 7.47 & 4.64 & 7.19 & 4.84 & \\
%\midrule[0.6pt]
%Equiformer + NAT & 0.4685 & 0.6071 & 0.4746 & 0.5345 & 0.5212 & 5.83 & 3.18 & 5.88 & 3.43 & \\
\bottomrule[1.2pt]
\end{tabular}
}
\end{adjustwidth}
\vspace{-3mm}
\caption{
%\revision{
%%%\textbf{Results on OC20 IS2RE \underline{testing} set when IS2RS node-level auxiliary task is adopted during training.} 
\textbf{Results on OC20 IS2RE \underline{testing} set when IS2RS is adopted during training.} 
$\dagger$ denotes using ensemble of models trained on both IS2RE training and validation sets.
In contrast, we use the same single Equiformer model in Table~\ref{tab:is2re_is2rs_validation}, which is trained only on the training set.
Note that Equiformer achieves better results with much less computation.
%}
}
\vspace{-6mm}
\label{tab:is2re_is2rs_testing}
\end{table}


\vspace{-3mm}
\subsection{OC20}
\vspace{-2mm}
\label{subsec:oc20}

\vspace{-1mm}
\paragraph{Dataset.}
%The Open Catalyst 2020 (OC20) dataset~\cite{oc20} contains 660k distinct large atomisitc systems, each consisting of a small molecule called adsorbate placed on a large slab called catalyst.
The Open Catalyst 2020 (OC20) dataset~\citep{oc20} (Creative Commons Attribution 4.0 License) consists of larger atomic systems, each composed of a molecule called adsorbate placed on a slab called catalyst.
Each input contains more atoms and more diverse atom types than QM9 and MD17.
%%%The average number of atoms in a system is more than 70, and there are over 50 atom species.
%%%The goal is to understand interaction between adsorbates and catalysts through relaxation.
%Specifically, an adsorbate is first placed on top of a catalyst at heuristically determined locations in order to form initial structure (IS).
%%%An adsorbate is first placed on top of a catalyst to form initial structure (IS).
%%%The positions of atoms are updated with forces calculated by density function theory until the system is stable and becomes relaxed structure (RS).
%%%The energy of RS, or relaxed energy (RE), is correlated with catalyst activity and therefore a metric for understanding their interaction.
%%%We focus on the task of initial structure to relaxed energy (IS2RE), which predicts relaxed energy (RE) given an initial structure (IS).
We focus on the task of initial structure to relaxed energy (IS2RE), which is to predict the energy of a relaxed structure (RS) given its initial structure (IS).
%%%There are 460k, 100k and 100k structures in training, validation, and testing sets, respectively.
Performance is measured in MAE and energy within threshold (EwT), the percentage in which predicted energy is within 0.02~eV of ground truth energy.
In validation and testing sets, there are four sub-splits containing in-distribution adsorbates and catalysts (ID), out-of-distribution adsorbates (OOD-Ads), out-of-distribution catalysts (OOD-Cat), and out-of-distribution adsorbates and catalysts (OOD-Both). 
%We minimize MAE between predictions and normalized ground truth.
Please refer to Sec.~\ref{appendix:subsec:oc20_dataset_details} for the detailed description of OC20 dataset.

\vspace{-3.25mm}
\paragraph{Setting.}
We consider two training settings based on whether a node-level auxiliary task~\citep{noisy_nodes} is adopted.
In the first setting, we minimize MAE between predicted energy and ground truth energy without any node-level auxiliary task.
%%%In the second setting, we incorporate the task of initial structure to relaxed structure (IS2RS) as a node-level auxiliary task. %, graphormer_3d}.
%%%In addition to predicting energy, we predict node-wise vectors indicating how each atom moves from initial structure to relaxed structure.
%%%In the second setting, we incorporate the task of initial structure to relaxed structure (IS2RS) as a node-level auxiliary task and additionally predict node-wise vectors indicating how each atom moves from initial structure to relaxed structure.
In the second setting, we incorporate the task of initial structure to relaxed structure (IS2RS) as a node-level auxiliary task.

\vspace{-3.25mm}
\paragraph{Training Details.}
Please refer to Sec.~\ref{appendix:subsec:oc20_training_details} in appendix for details on Equiformer architecture, hyper-parameters and training time.


\vspace{-3.25mm}
\paragraph{IS2RE Results without Node-Level Auxiliary Task.}
We summarize the results on validation and testing sets under the first setting in Table~\ref{tab:is2re_validation} in appendix and Table~\ref{tab:is2re_test}.
Compared with state-of-the-art models like SEGNN and SphereNet, Equiformer consistently achieves the lowest MAE for all the four sub-splits in validation and testing sets.
%%%Note that energy within threshold (EwT) considers only the percentage of predictions close enough to ground truth and the distribution of errors, and therefore improvement in average errors (MAE) would not necessarily reflect that in error distributions (EwT).
Note that EwT considers only the percentage of predictions close enough to ground truth and the distribution of errors, and therefore improvement in average errors (MAE) would not necessarily reflect that in error distributions (EwT).
A more detailed discussion can be found in Sec.~\ref{appendix:subsec:error_distributions} in appendix.
%as long as errors are still higher than 0.02~eV.
%Similar phenomena can be observed in OC20 leaderboard\footnote{As shown in the leaderboard (\href{https://opencatalystproject.org/leaderboard.html}{https://opencatalystproject.org/leaderboard.html}), ``GNS + Noisy Nodes (Direct)'' achieves lower MAE but lower EwT than ``DimeNet++-1.8M-fonly-eonly-relaxation-All''. 
%A screenshot of the leaderboard can be found \href{https://drive.google.com/file/d/1nxS_Zb4VrV1nIEAhkPFaj9mcqaHfQiL_/view?usp=sharing}{here}.}.
%%%Similar phenomena can be observed in Table~\ref{tab:is2re_test}, where for ``OOD Both'' sub-split, SphereNet achieves lower MAE yet lower EwT than SEGNN.
%%%We also note that models in Table~\ref{tab:is2re_validation} and~\ref{tab:is2re_test} are trained by minimizing MAE, and therefore comparing MAE in validation and testing sets could mitigate the discrepancy between training objectives and evaluation metrics and that OC20 leaderboard ranks the relative performance of models mainly according to MAE.
We also note that models are trained by minimizing MAE, and therefore comparing MAE could mitigate the discrepancy between training objectives and evaluation metrics and that OC20 leaderboard ranks the relative performance according to MAE.
Additionally, we compare the training time of SEGNN and Equiformer in Sec.~\ref{appendix:subsec:oc20_training_time_comparison}.

\vspace{-3.5mm}
\paragraph{IS2RE Results with IS2RS Node-Level Auxiliary Task.}
We report the results on validation and testing sets in Table~\ref{tab:is2re_is2rs_validation} and Table~\ref{tab:is2re_is2rs_testing}.
As of the date of the submission of this work, Equiformer achieves the best results on IS2RE task when only IS2RE and IS2RS data are used.
%%%Notably, the proposed Equiformer in Table~\ref{tab:is2re_is2rs_testing} achieves competitive results even with much less computation.
Notably, the result in Table~\ref{tab:is2re_is2rs_testing} is achieved with much less computation.
%Specifically, training ``Equiformer + Noisy Nodes'' takes about $24$ GPU-days when A6000 GPUs are used.
%The training time of ``GNS + Noisy Nodes'' is $56$ TPU-days.
%``Graphormer'' uses ensemble of $31$ models and requires $372$ GPU-days when A100 GPUs are used.
%The comparison to GNS demonstrates the improvement from invariant message passing networks to equivariant Transformers.
%Without any data augmentation, Equiformer still achieves competitive results to GNS trained with Noisy Nodes~\cite{noisy_nodes}.
%Compared to Graphormer, Equiformer demonstrates the effectiveness of equivariant features and the proposed equivariant graph attention.
We note that under this setting, greater depths and thus more computation translate to better performance~\citep{noisy_nodes} and that Equiformer demonstrates incorporating equivariant features and the proposed equivariant graph attention can improve training efficiency by $2.3\times$ to $15.5\times$ compared to invariant message passing networks and invariant Transformers.
%%%Note that Equiformer, with $18$ Transformer blocks, is relatively shallow as GNS trained with Noisy Nodes has $100$ blocks and Graphormer has $48$ Transformer blocks and that deeper networks can typically obtain better results when IS2RS auxiliary task is adopted~\cite{noisy_nodes}.
